% This document is compiled using pdfLaTeX
% You can switch XeLaTeX/pdfLaTeX/LaTeX/LuaLaTeX in Settings

\documentclass{article}
\usepackage{ctex}
\usepackage{amsmath} % 确保引入了此宏包以支持 aligned

\title{因子日历2026}

\date{\today}

\begin{document}

\maketitle

\section{资本利得突出量(行为金融因子-处置效应)}

\subsection{因子定义}
使用过去一段时间的换手率加权价格作为参考价格，用于衡量股票的综合持仓成本，进而衡量大部分投资者在某只股票上是盈利还是亏损：
\begin{equation}
    \mathrm{RP}_t = \frac{1}{k} \sum_{n=0}^{T-1} \left( V_{t-n} \prod_{s=1}^{n-1} (1 - V_{t-n+s}) \right) P_{t-n}
\end{equation}
计算第 \( t \) 日股票 \( i \) 的持仓者相对参考价格数学 \( \mathrm{RP}_i \) 的平均盈亏情况即为资本利得突出量：
\begin{equation}
    \mathrm{CGO}_{i,t} = \frac{\mathrm{Close}_{i,t} - \mathrm{RP}_t}{\mathrm{Close}_{i,t}}
\end{equation}

\subsection{变量定义}
\begin{itemize}
    \item \( V_t \)、\( \mathrm{Close} \)、\( P_t \) 分别为为 \( t \) 日的换手率、收盘价、VWAP 均价；
    \item \( k \) 为权重系数，用于保证计算 \( \mathrm{RP}_t \) 时的权重和等于 1，即为 \( P_t \) 前面权重系数之和。
\end{itemize}

\subsection{说明}
处置效应指投资者“售盈持亏”的情感倾向；CGO 因子值越大，浮盈越大；历史股价上涨，高 CGO 意味着大量投资者盈利，抛压大，卖盘压力增加，股价上涨受阻；历史股价下跌，低或负 CGO 意味着投资者普遍被套，抛售意愿弱，卖压小，一旦有利好，股价更容易反弹；CGO 因子对应反转效应，与未来收益负相关。

\subsection{参考文献}
\begin{enumerate}
    \item 陈开锐, 姚紫薇. 2025, 处置效应与V型处置效应对在量化选股中的应用. 中信建投.
    \item Grinblatt, M., \& Han, B. (2005). \textit{Prospect theory, mental accounting, and the disposition effect}. Journal of Financial Economics, 78(2), 311–339.
\end{enumerate}

\section{残差反转}

\subsection{因子定义（高频因子-动量反转类）}

计算小单资金流强度：分子为（小单买额-小单卖额）之和，分母为其绝对值；计算大单资金流强度：分子为（小单买额-小单卖额）之和。
\begin{equation}
    \mathrm{S}_t = \frac{\sum_{t-T}^t (\mathrm{buy}_t - \mathrm{sell}_t)}{\sum_{t-T}^t |\mathrm{buy}_t - \mathrm{sell}_t|}
\end{equation}
将反转因子（过去20日涨跌幅）与大单资金强度和小单资金强度各自做回归，得到的残差即为得到的大单和小单的残差反转因子：
\begin{equation}
    \mathrm{Ret20}_t = a + b * \mathrm{S}_t + \varepsilon_t
\end{equation}

\subsection{说明}
残差反转因子剥离了资金流强度对传统反转因子的影响，选股效果得到进一步提升。

\subsection{参考文献}
\begin{enumerate}
    \item 张建超, 高鹏. 2021. 大单与小单资金流的 alpha 能力. 开源证券.
\end{enumerate}

\section{个人投资者交易热度（行为金融因子-投资者注意力）}

\subsection{因子定义}

\begin{equation}
    \mathrm{TURN\_RETAIL} = \frac{1}{m} \sum_{t=1}^{t-m+1} \frac{\mathrm{INFLOW\_RETAIL}_{i,t} + \mathrm{OUTFLOW\_RETAIL}_{i,t}}{\mathrm{CAP}_{i,t}}
\end{equation}

\subsection{变量定义}
\begin{itemize}
    \item ① \(\mathrm{INFLOW\_RETAIL}_{i,t}\) 和 \(\mathrm{OUTFLOW\_RETAIL}_{i,t}\) 分别为股票 \(i\) 在交易日 \(t\) 的个人投资者买入金额和卖出金额；
    \item ② 将单笔成交额低于 4 万元的投资者作为个人投资者；
    \item ③ \(\mathrm{CAP}_{i,t}\) 为股票 \(i\) 在交易日 \(t\) 的流通市值；
    \item ④ \(m\) 为交易日窗口，一般取为月度。
\end{itemize}

\subsection{说明}
个人投资者交易热度属于股票交易活跃度类因子，通过衡量股票交易活跃度来反映投资者关注度。股票交易越活跃，投资者对该股票的关注度越高；而投资者有限注意力会推动关注度高的股票更易出现非理性净买入溢价，进而导致股价后续反转。

\subsection{参考文献}
\begin{enumerate}
    \item 陈升锐. 2024. 投资者有限关注及注意力捕捉与溢出. 中信建投.
\end{enumerate}

\section{换手率分布均匀度(高频因子-成交分布类)}

\subsection{因子定义}


每月月底，回溯所有股票过去20个交易日，计算当日分钟换手率的标准差：
\begin{equation}
    \mathrm{TurnVolDaily} = \mathrm{std}(\text{分钟换手率})
\end{equation}
其中，
\begin{equation}
    \text{分钟换手率} = \frac{\text{分钟成交量}}{\text{当日流通股本}}
\end{equation}
基于这20个交易日的分钟换手率序列，计算标准差和均值，然后进行标准化和市值中性化处理得到换手率分布均匀度 \(\mathrm{UTD}\)：
\begin{equation}
    \mathrm{UTD} = \frac{\mathrm{std}(\mathrm{TurnVolDaily})}{\mathrm{mean}(\mathrm{TurnVolDaily})}
\end{equation}

\subsection{说明}
换手率分布均匀度因子衡量了股票换手率在日内不同交易时段的分布均匀程度，以及该均匀程度在日间是否平稳，股票换手率在日内及日间越平稳，即因子值越小越好。

\subsection{参考文献}
\begin{enumerate}
    \item 高子剑. 2021. 换手率分布均匀度：基于分钟成交量的选股因子. 东吴证券.
\end{enumerate}

\section{筹码分布成本重心(行为金融因子-筹码分布)}

\subsection{因子定义}

\begin{equation}
    \mathrm{CKDP}_t = \frac{\mathrm{mean}_t - \mathrm{Chip\_Lowest}_t}{\mathrm{Chip\_Highest}_t - \mathrm{Chip\_Lowest}_t}
\end{equation}

\subsection{变量定义}
\begin{itemize}
    \item ① 假设某只股票在 \( t \) 日的筹码分布共有 \( n \) 个价位（\( i = 1, 2, ..., n \)），\( \mathrm{vol}_{i,t} \) 为第 \( t \) 日 \( i \) 价位上的筹码峰高度，\( \mathrm{price}_{i,t} \) 为该筹码峰对应的价格，\( \mathrm{Chip\_Highest}_t \)、\( \mathrm{Chip\_Lowest}_t \) 分别对应筹码最高价和筹码最低价；
    \item ② \(\displaystyle \mathrm{mean}_t = \sum (\mathrm{price}_{i,t} * p_{i,t})\) 为筹码分布加权平均成本。
\end{itemize}

\subsection{说明}
筹码分布成本重心衡量了筹码平均成本在筹码分布中的相对位置，可用于识别建仓或出货阶段。成本重心接近 0，说明平均成本靠近成本区间下沿位置，成本偏低，可能处于低位建仓阶段；成本重心接近 1，说明平均成本靠近区间上沿位置，成本偏高，可能处于高位派发阶段；成本重心接近 0.5，成本分布对称，市场处于平衡状态。经测试，筹码分布成本宽带与未来收益负相关。

\subsection{参考文献}
\begin{enumerate}
    \item 陈升锐, 姚紫薇. 2025. 筹码分布因子系统构建. 中信建投.
\end{enumerate}

\section{下行跳跃波动率(高级因子-波动跳跃类)}

\subsection{因子定义}

\begin{equation}
    \mathrm{RVJN}_t = \max\left(\mathrm{RS}_t^- - \frac{1}{2} \hat{IV}_t, 0\right)
\end{equation}
\begin{equation}
    \mathrm{RS}_t^- = \sum_{i=1}^n r_{t,i}^2 \mathrm{I}_{r_{t,i} > 0} \rightarrow_u \frac{1}{2} \int_0^t \sigma_s^2 ds + \sum_{s \leq t} \Delta X_s^2 \mathrm{I}_{\Delta X_s < 0}
\end{equation}
\begin{equation}
    \hat{IV}_t = \mu^{-3}_{\frac{2}{3}} \sum_{i=k}^n |r_{t_i}|^{\frac{2}{3}} |r_{t_{i-1}}|^{\frac{2}{3}} \cdots |r_{t_{i-k+1}}|^{\frac{2}{3}}
\end{equation}

\subsection{变量定义}
\begin{itemize}
    \item ① \( r_{t_i} \) 为某股票交易日 \( t \) 内的分钟对数收益率序列，如 5 分钟对数收益率序列；
    \item ② \( k=3 \)，\( n \) 表示该股票在交易日 \( t \) 中特定频率的分钟级别数据个数，如 5 分钟频率通常有 48 个数据点；
    \item ③ \( \mu_q \) 为标准正态分布的第q阶矩对短；
    \item ④ 下行跳跃波动率 \(\mathrm{RVJN}\) 整体取值为正，所以需要用 \(\max\) 对负差值做修正。
\end{itemize}

\subsection{说明}
按收益涨跌情况，股价波动可分为上行波动和下行波动，跳跃波动也可据此做划分。下行跳跃波动就是下行波动中的跳跃部分；短期内，下行跳跃波动通常有正的风险溢价。

\subsection{参考文献}
\begin{enumerate}
    \item Bo Yu, Bruce Mizrach, Norman R. Swanson. (2020). \textit{New Evidence of the Marginal Predictive Content of Small and Large Jumps in the Cross-Section}. Econometrics, MDPI, 8(2), 1-52.
\end{enumerate}

\section{下游传导动量(图谱网络-动量溢出)}

\subsection{因子定义}
计算上游公司 A 与下游公司 B 间每一对上下游产品 \(i \to j\) 之间的关联度：
\begin{equation}
    \mathrm{wgt}_{A,B,t} = \frac{\sum_i \mathrm{prdIncome}_{A,i,t} \times \mathbf{1}\{i \to j, i \in A\} + \sum_j \mathrm{prdIncome}_{B,j,t} \times \mathbf{1}\{i \to j, j \in B\}}{\sum_i \mathrm{prdIncome}_{A,i,t} + \sum_j \mathrm{prdIncome}_{B,j,t}}
\end{equation}
对上游公司 A，以产业链关联度为权重加权其所有下游公司的月度收益率，即为下游传导动量因子：
\begin{equation}
    \mathrm{upstream\_Moment}_{A,t} = \frac{\sum_{A \neq B} \mathrm{wgt}_{A,B,t} \times \mathrm{rtn}_{B,t}}{\sum_{A \neq B} \mathrm{wgt}_{A,B,t}}
\end{equation}

\subsection{变量定义}
\begin{itemize}
    \item ① \(\mathrm{prdIncome}\) 为公司 A 在产品 i 上的营业收入；
    \item ② \(\sum_i \mathrm{prdIncome}_{A,i,t}\) 为公司 A 的当期总营收；
    \item ③ \(\mathbf{1}\{i \to j\}\) 表示产品 i 为产品 j 的上游；
    \item ④ \(\mathrm{rtn}_{B,t}\) 为公司 B 在 t 时刻的月度收益率；
    \item ⑤ 可以将下游传导动量因子对月度收益率做正交化，剥离月度反转效应。
\end{itemize}

\subsection{说明}
下游传导动量因子值越大，表示公司 A 作为上游企业时，位于其产业链下游的高关联度高集群公司在考察月份涨幅相对较高，公司 A 接下来也会有上涨的可能。

\subsection{参考文献}
\begin{enumerate}
    \item 郑兆磊.2023.产业链视角下的Alpha传导研究.兴业证券
\end{enumerate}

\section{超大单和小单的同步相关性（高频因子-资金流类）}

\subsection{因子定义}

\begin{equation}
    \mathrm{RankCorr}(\mathrm{EL}_t, \mathrm{S}_t)
\end{equation}

\subsection{变量定义}
\begin{itemize}
    \item ① \( \mathrm{EL}_t \) 为超大单（>100万元）过去20个交易日的资金净流入序列；
    \item ② \( \mathrm{S}_t \) 为小单（<4万元）过去20个交易日同期的资金净流入序列。
\end{itemize}

\subsection{说明}
在资金流同步相关性因子中，上述超大单和小单的同步相关性因子表现最好；但资金流同步相关性因子是“伪动力学”因子，在恒等式“超大单+大单+中单+小单=0”的约束下，该因子的 alpha，可以被两者流入量级所解释，其底层 alpha 来源仍是资金流强度。

\subsection{参考文献}
\begin{enumerate}
    \item 魏建榕, 盛少成. 2022. 新型因子：资金流动力学与散户羊群效应. 开源证券.
\end{enumerate}

\section{买入浮盈系数修正后的积极灵活筹码（行为金融因子-筹码分布）}

\subsection{因子定义}
计算筹码当日收盘价相对筹码买入价格的涨跌幅作为该筹码的买入浮盈系数：
\begin{equation}
    a_i = \frac{\mathrm{Close}_t}{\mathrm{TradePrice}_i}
\end{equation}
将买入浮盈系数与筹码买入量相乘即可得到修正后的积极灵活筹码值：
\begin{equation}
    \mathrm{AFH\_Close} = \frac{\sum_{i=1}^N a_i * \mathrm{Vol\_AFH}_i}{\sum_{t=T-20}^T \mathrm{Volume}_t}
\end{equation}

\subsection{变量定义}
\begin{itemize}
    \item ① \(i\) 为筹码购入时的订单号，统计对象为过去 20 日的订单号；
    \item ② \( \mathrm{Vol\_AFH}_i \) 为主动买入时订单金额小于 5 万元的订单成交量；
    \item ③ \( \mathrm{TradePrice}_i \) 为该订单的买入价格；
    \item ④ \( \mathrm{Close}_t \)、\( \mathrm{Volume}_t \) 分别为订单买入当日的开盘价和当日总成交量。
\end{itemize}

\subsection{说明}
投资者会由自我归因偏差引起过度自信，这将加强投资者的处置效应。当投资者买入相应个股之后，若当日即获得浮盈，将使投资者形成“买对了”的心理行为，这将提升投资者对股票判断的自信，此时的买入筹码也更具处置效应，因此可以计算筹码当日收盘价相对筹码买入价格的涨跌幅作为该筹码的买入浮盈系数，用于修正原始的积极灵活筹码，提升因子表现。

\subsection{参考文献}
\begin{enumerate}
    \item 任瞳, 许继宏. 2024. 基于逐笔数据的 AFH 改进因子. 招商证券.
\end{enumerate}

\section{非流动性（高频因子-流动性类）}

\subsection{因子定义}

不考虑日内价格变动方向，直接进行无向叠加，求得股价变动的最短路径：
\begin{equation}
    \mathrm{shortcut} = 2 * (\mathrm{high} - \mathrm{low}) - |\mathrm{close} - \mathrm{open}|
\end{equation}
利用高频分钟 K 线，每日的 K 线最短路径非流动性因子：
\begin{equation}
    \mathrm{ILLIQ}_t = \frac{1}{d} \sum_{i=1}^d \left( \sum_{j=1}^p \frac{\mathrm{ShortCut}_j}{\mathrm{Value}_j} \right)_{t-i}
\end{equation}

\subsection{变量定义}
\begin{itemize}
    \item ① \(\mathrm{ShortCut}_j\) 表示日内频率下第 \(j\) 根 K 线的最短路径；
    \item ② \(\mathrm{Value}_j\) 表示日内频率下第 \(j\) 根 K 线内的成交额；
    \item ③ \(p\) 表示日内频率分段个数；
    \item ④ \(d\) 表示日间移动平均的周期参数。
\end{itemize}

\subsection{说明}
通过使用更高频的 K 线数据提升 K 线最短路径对股票交易时市场冲击的代理精度，大大提升了 K 线最短路径非流动性因子的有效性与预测能力。同时随着数据频率的提高，K 线最短路径定义对经典定义下的非流动性因子提升程度愈加显著。比起经典的经典定义方式，K 线最短路径能更充分地利用高频数据新引入的信息。

\subsection{参考文献}
\begin{enumerate}
    \item 刘均伟. 2017. 基于 K 线最短路径构造的非流动性因子. 光大证券.
\end{enumerate}
\end{document}